\scanpage{47}\mspage{46}
Un ensemble \uline{complet} de conditions \unsure{de} compatibilité-\unsure{disjonction|des parties} impliquant $\Lambda$ \deleted{\unsure{peut}} (impliquant donc tous les axiomes pertinents
\diagram{\begin{tikzcd}[ampersand replacement=\&,column sep=.22em,row sep=.55em]
C, \& D,\arrow[d,Rightarrow] \& C\mathord{*},\ \Lambda1,\ \Lambda2,\ \Lambda3,\ C\Lambda\mathord{*}\\
{} \& D_0 \& {}
\end{tikzcd}}
peut être caractérisé de diverses façons, p. ex

\diagram{\begin{tikzcd}[ampersand replacement=\&,column sep=.75em,row sep=.8em,execute at end picture={
\path (b.west) -- (c.west) coordinate[midway] (mid);
\draw ([xshift=-2pt]a.north west) .. controls +(-4pt,0pt) and +(0pt,3pt) .. ([xshift=-5pt]a.west) -- ([xshift=-5pt,yshift=4pt]mid) .. controls +(0pt,-3pt) and +(3pt,0pt) .. ([xshift=-9pt]mid) .. controls +(3pt,0pt) and +(0pt,3pt) .. ([xshift=-5pt,yshift=-4pt]mid) -- ([xshift=-5pt]d.west) .. controls +(0pt,-3pt) and +(-4pt,0pt) .. ([xshift=-2pt]d.south west);}]
a) \& |[alias=a]|\Lambda1, \& \Lambda2, \& \Lambda3,\arrow[d,phantom,"\unsure{\Updownarrow|\Downarrow}" description] \& C\Lambda\mathord{*}\arrow[d,Leftrightarrow]\\
b) \& |[alias=b]|\Lambda1, \& \Lambda2, \& D, \& C\mathord{*}\\
c) \& |[alias=c]|\Lambda1, \& C\Lambda\mathord{*}, \& D \& {}\\
d) \& |[alias=d]|\Lambda1,\unsure{''} \& \deleted{\unsure{C\mathord{*}}}\Lambda2, \& \deleted{\unsure{D_0}}C\mathord{*}, \& D_0
\end{tikzcd}}
\margin{les équivalences marquées \unsure{valent|valant} terme à terme, en présence de $\Lambda1,\ \Lambda2$ i.e. si $\Lambda$ est partie très dense de $M$}
\margin{\uline{NB} $\Lambda1$ est le seul des axiomes partiels qui figure dans \unsure{chacun} des quatre \unsure{syst.} équivalents a) b) c) d)}

On dira, quand ces conditions sont vérifiées qu'on a un atlas $\Lambda$-spacieux, ou $\Lambda$-polyédral. Les cas les plus intéressants pour moi sont

$1^{\circ})$ $\Lambda=L$ cas déjà traité --- mais cela suppose que l'atlas satisfait notamment la condition de divisibilité

$2^{\circ})$ $\Lambda=M$, auquel cas $\Lambda1$ \deleted{\unsure{est}} devient tautologique, $\Lambda2$ équivaut à $D_0$, $\Lambda3$ équivaut à $D$, $C\Lambda\mathord{*}$ équivaut à $C\mathord{*}$, \deleted{et la} Dire que $A$ est $M$-spacieux \inserted{ou $M$-polyédr\unsure{.}} signifie donc simplement qu'il satisfait \unsure{:|à} $D$, et à $C$ spac \unsure{resp} $C$ pol.

Dans le cas général, on peut dire que $A$ est (disons) $\Lambda$-spacieux, \unsure{dès|ssi|si} il est $M$-spacieux i.e. spacieux tout court, et si de plus $\Lambda$ est très dense ($\Lambda1$ et $\Lambda2$)
